
\exo

\vspace{-12pt}
\setlength{\columnsep}{1cm}
\setlength{\columnseprule}{0pt}
\begin{multicols}{2}
\begin{enumerate}
\item Déterminer les coordonnées des vecteurs $\vecu$, $\vecv$, $\vecw$ et $\vect{t}$ dans la base $\baseij$.

\begin{center}
\def\xmin{0}
\def\xmax{13}
\def\ymin{0}
\def\ymax{5}
\def\xunite{0.6cm}
\def\yunite{0.6cm}
\psset{unit=\xunite, xunit=\xunite, yunit=\yunite, algebraic=true, dimen=middle, dotstyle=*, dotsize=3pt 0, linewidth=0.8pt, arrowsize=3pt 2, arrowinset=0.25, comma=true}
\begin{pspicture*}(\xmin,\ymin)(\xmax,\ymax)
\psgrid[xunit=\xunite, yunit=\yunite, subgriddiv=0, gridlabels=0, gridcolor=lightgray](0,0)(\xmin,\ymin)(\xmax,\ymax)
%\psaxes[labelFontSize=\scriptstyle, xAxis=true, yAxis=true, Dx=1, Dy=1, ticksize=-2pt 0, subticks=1]{->}(0,0)(\xmin,\ymin)(\xmax,\ymax)[{\footnotesize $x$},135][{\footnotesize $y$},-45]
\psline[linewidth=1.2pt,arrows=->](5,1)(6,1)
\psline[linewidth=1.2pt,arrows=->](5,1)(5,2)
\psline[arrows=->](4,1)(2,2)
\psline[arrows=->](4,4)(1,2)
\psline[arrows=->](4,3)(7,3)
\psline[arrows=->](7,1)(11,3)
\begin{footnotesize}
\uput{3pt}[270](5.5,1){$\veci$}
\uput{3pt}[180](5,1.5){$\vecj$}
\uput{3pt}[060](3,1.5){$\vecu$}
\uput{3pt}[120](2.5,3){$\vecv$}
\uput{3pt}[120](9,2){$\vecw$}
\uput{3pt}[090](5.5,3){$\vect{t}$}
\end{footnotesize}
\end{pspicture*}
\end{center}

\item Déterminer les coordonnées des vecteurs $\vecu$, $\vecv$, $\vecw$ et $\vect{t}$ dans la base $\baseij$.

\begin{center}
\def\xmin{0}
\def\xmax{13}
\def\ymin{0}
\def\ymax{5}
\def\xunite{0.6cm}
\def\yunite{0.6cm}
\psset{unit=\xunite, xunit=\xunite, yunit=\yunite, algebraic=true, dimen=middle, dotstyle=*, dotsize=3pt 0, linewidth=0.8pt, arrowsize=3pt 2, arrowinset=0.25, comma=true}
\begin{pspicture*}(\xmin,\ymin)(\xmax,\ymax)
\psgrid[xunit=\xunite, yunit=\yunite, subgriddiv=0, gridlabels=0, gridcolor=lightgray](0,0)(\xmin,\ymin)(\xmax,\ymax)
%\psaxes[labelFontSize=\scriptstyle, xAxis=true, yAxis=true, Dx=1, Dy=1, ticksize=-2pt 0, subticks=1]{->}(0,0)(\xmin,\ymin)(\xmax,\ymax)[{\footnotesize $x$},135][{\footnotesize $y$},-45]
\psline[linewidth=1.2pt,arrows=->](1,1)(3,1)
\psline[linewidth=1.2pt,arrows=->](1,1)(1,3)
\psline[arrows=->](4,1)(4,4)
\psline[arrows=->](5,1)(9,4)
\psline[arrows=->](11,1)(7,2)
\psline[arrows=->](12,3)(9,3)
\begin{footnotesize}
\uput{3pt}[270](2,1){$\veci$}
\uput{3pt}[180](1,2){$\vecj$}
\uput{3pt}[130](7,2.5){$\vecu$}
\uput{3pt}[180](4,2.5){$\vecv$}
\uput{3pt}[060](9,1.5){$\vecw$}
\uput{3pt}[090](10.5,3){$\vect{t}$}
\end{footnotesize}
\end{pspicture*}
\end{center}
\end{enumerate}
\end{multicols}
\vspace{-12pt}

\exo

Dans la base $\baseij$ de chacune des figures suivantes, représenter les vecteurs :
\[\vecu\couple{3}{-2}
\qpvq
\vecv\couple{-1}{2}
\qpvq
\vecw\couple{0}{-3}
\qpvq
\vect{a}\couple{3}{0}.
\]

\vspace{-12pt}
\setlength{\columnsep}{1cm}
\setlength{\columnseprule}{0pt}
\begin{multicols}{2}
\begin{center}
\def\xmin{0}
\def\xmax{13}
\def\ymin{0}
\def\ymax{5}
\def\xunite{0.6cm}
\def\yunite{0.6cm}
\psset{unit=\xunite, xunit=\xunite, yunit=\yunite, algebraic=true, dimen=middle, dotstyle=*, dotsize=3pt 0, linewidth=0.8pt, arrowsize=3pt 2, arrowinset=0.25, comma=true}
\begin{pspicture*}(\xmin,\ymin)(\xmax,\ymax)
\psgrid[xunit=\xunite, yunit=\yunite, subgriddiv=0, gridlabels=0, gridcolor=lightgray](0,0)(\xmin,\ymin)(\xmax,\ymax)
%\psaxes[labelFontSize=\scriptstyle, xAxis=true, yAxis=true, Dx=1, Dy=1, ticksize=-2pt 0, subticks=1]{->}(0,0)(\xmin,\ymin)(\xmax,\ymax)[{\footnotesize $x$},135][{\footnotesize $y$},-45]
\psline[linewidth=1.2pt,arrows=->](1,1)(2,1)
\psline[linewidth=1.2pt,arrows=->](1,1)(1,2)
\begin{footnotesize}
\uput{3pt}[270](1.5,1){$\veci$}
\uput{3pt}[180](1,1.5){$\vecj$}
\end{footnotesize}
\end{pspicture*}
\end{center}

\begin{center}
\def\xmin{0}
\def\xmax{13}
\def\ymin{0}
\def\ymax{5}
\def\xunite{0.6cm}
\def\yunite{0.6cm}
\psset{unit=\xunite, xunit=\xunite, yunit=\yunite, algebraic=true, dimen=middle, dotstyle=*, dotsize=3pt 0, linewidth=0.8pt, arrowsize=3pt 2, arrowinset=0.25, comma=true}
\begin{pspicture*}(\xmin,\ymin)(\xmax,\ymax)
\psgrid[xunit=\xunite, yunit=\yunite, subgriddiv=0, gridlabels=0, gridcolor=lightgray](0,0)(\xmin,\ymin)(\xmax,\ymax)
%\psaxes[labelFontSize=\scriptstyle, xAxis=true, yAxis=true, Dx=1, Dy=1, ticksize=-2pt 0, subticks=1]{->}(0,0)(\xmin,\ymin)(\xmax,\ymax)[{\footnotesize $x$},135][{\footnotesize $y$},-45]
\psline[linewidth=1.2pt,arrows=->](1,1)(3,1)
\psline[linewidth=1.2pt,arrows=->](1,1)(1,2)
\begin{footnotesize}
\uput{3pt}[270](2,1){$\veci$}
\uput{3pt}[180](1,1.5){$\vecj$}
\end{footnotesize}
\end{pspicture*}
\end{center}

\begin{center}
\def\xmin{0}
\def\xmax{13}
\def\ymin{0}
\def\ymax{5}
\def\xunite{0.6cm}
\def\yunite{0.6cm}
\psset{unit=\xunite, xunit=\xunite, yunit=\yunite, algebraic=true, dimen=middle, dotstyle=*, dotsize=3pt 0, linewidth=0.8pt, arrowsize=3pt 2, arrowinset=0.25, comma=true}
\begin{pspicture*}(\xmin,\ymin)(\xmax,\ymax)
%\psgrid[xunit=\xunite, yunit=\yunite, subgriddiv=0, gridlabels=0, gridcolor=lightgray](0,0)(\xmin,\ymin)(\xmax,\ymax)
%\psaxes[labelFontSize=\scriptstyle, xAxis=true, yAxis=true, Dx=1, Dy=1, ticksize=-2pt 0, subticks=1]{->}(0,0)(\xmin,\ymin)(\xmax,\ymax)[{\footnotesize $x$},135][{\footnotesize $y$},-45]
\psline[linecolor=lightgray](\xmin,0)(\xmax,0)
\psline[linecolor=lightgray](\xmin,1)(\xmax,1)
\psline[linecolor=lightgray](\xmin,2)(\xmax,2)
\psline[linecolor=lightgray](\xmin,3)(\xmax,3)
\psline[linecolor=lightgray](\xmin,4)(\xmax,4)
\psline[linecolor=lightgray](\xmin,5)(\xmax,5)
\psplot[linecolor=lightgray]{\xmin}{\xmax}{2*x+4}
\psplot[linecolor=lightgray]{\xmin}{\xmax}{2*x+2}
\psplot[linecolor=lightgray]{\xmin}{\xmax}{2*x+0}
\psplot[linecolor=lightgray]{\xmin}{\xmax}{2*x-2}
\psplot[linecolor=lightgray]{\xmin}{\xmax}{2*x-4}
\psplot[linecolor=lightgray]{\xmin}{\xmax}{2*x-6}
\psplot[linecolor=lightgray]{\xmin}{\xmax}{2*x-8}
\psplot[linecolor=lightgray]{\xmin}{\xmax}{2*x-10}
\psplot[linecolor=lightgray]{\xmin}{\xmax}{2*x-12}
\psplot[linecolor=lightgray]{\xmin}{\xmax}{2*x-14}
\psplot[linecolor=lightgray]{\xmin}{\xmax}{2*x-16}
\psplot[linecolor=lightgray]{\xmin}{\xmax}{2*x-18}
\psplot[linecolor=lightgray]{\xmin}{\xmax}{2*x-20}
\psplot[linecolor=lightgray]{\xmin}{\xmax}{2*x-22}
\psplot[linecolor=lightgray]{\xmin}{\xmax}{2*x-24}
\psplot[linecolor=lightgray]{\xmin}{\xmax}{2*x-26}
\psline[linewidth=1.2pt,arrows=->](1.5,1)(3.5,1)
\psline[linewidth=1.2pt,arrows=->](1.5,1)(2,2)
\begin{footnotesize}
\uput{3pt}[270](2.5,1){$\veci$}
\uput{5pt}[180](1.75,1.5){$\vecj$}
\end{footnotesize}
\end{pspicture*}
\end{center}

\begin{center}
\def\xmin{0}
\def\xmax{13}
\def\ymin{0}
\def\ymax{5}
\def\xunite{0.6cm}
\def\yunite{0.6cm}
\psset{unit=\xunite, xunit=\xunite, yunit=\yunite, algebraic=true, dimen=middle, dotstyle=*, dotsize=3pt 0, linewidth=0.8pt, arrowsize=3pt 2, arrowinset=0.25, comma=true}
\begin{pspicture*}(\xmin,\ymin)(\xmax,\ymax)
\psgrid[xunit=\xunite, yunit=\yunite, subgriddiv=0, gridlabels=0, gridcolor=lightgray](0,0)(\xmin,\ymin)(\xmax,\ymax)
%\psaxes[labelFontSize=\scriptstyle, xAxis=true, yAxis=true, Dx=1, Dy=1, ticksize=-2pt 0, subticks=1]{->}(0,0)(\xmin,\ymin)(\xmax,\ymax)[{\footnotesize $x$},135][{\footnotesize $y$},-45]
\psline[linewidth=1.2pt,arrows=->](1,1)(3,1)
\psline[linewidth=1.2pt,arrows=->](1,1)(2,2)
\begin{footnotesize}
\uput{3pt}[270](2,1){$\veci$}
\uput{3pt}[180](1,1.5){$\vecj$}
\end{footnotesize}
\end{pspicture*}
\end{center}
\end{multicols}

\exo

Soient les vecteurs $\vecu\couple{5}{1}$ et $\vecv\couple{2}{-3}$. Calculer les coordonnées des vecteurs :
\[\vecu+\vecv
\qpvq
2\vecu
\qpvq
-3\vecv
\qpvq
2\vecu-3\vecv.\]

\exo

Soient les vecteurs $\vecu\couple{-2}{1}$ et $\vecv\couple{5}{-1}$. Calculer les coordonnées des vecteurs :
\[
\vecu+\vecv
\qpvq
\vecu-\vecv
\qpvq
3\vecu+4\vecv
\qpvq
-5\vecu+2\vecv.\]

\exo

Soient les vecteurs $\vecu\couple{3}{4}$ et $\vecv\couple{5}{-2}$.

\begin{enumerate}
\item Calculer la norme des vecteurs $\vecu$ et $\vecv$.
\item Calculer les normes :
\[
\norme{3\vecu}
\qpvq
\norme{\tfrac{1}{2}\vecu}
\qpvq
\norme{-2\vecv}
\qpvq
\norme{\vecu+\vecv}.
\]
\end{enumerate}
\exo

Dans chaque cas, placer un point $A$ dans un repère puis construire le point $M$ tel que $\vect{AM}=\vecu$.
\vspace{-9pt}
\setlength{\columnsep}{1cm}
\setlength{\columnseprule}{0pt}
\begin{multicols}{4}
\begin{enumerate}
\item $\vecu\couple{2}{1}$
\item $\vecu\couple{-3}{2}$
\item $\vecu\couple{0}{3}$
\item $\vecu\couple{2}{0}$
\end{enumerate}
\end{multicols}

\exo

Dans un repère, on considère les points :
\[
M\couple{-3}{6}
\qpvq
N\couple{5}{-1}
\qpvq
P\couple{11}{0}
\qpvq
Q\couple{3}{7}
\qpvq
R\couple{-10}{5}.
\]
\begin{enumerate}
\item Montrer que les vecteurs $\vect{MN}$ et $\vect{QP}$ sont égaux. Que peut-on en déduire ?
\item Le quadrilatère $MPNR$ est-il un parallélogramme ? Justifier la réponse.
\end{enumerate}

\exo

Dans un repère, on considère les points :
\[
E\couple{-1}{7}
\qpvq
F\couple{3}{-1}
\qpvq
G\couple{7}{-9}
\qpvq
H\couple{11}{-17}.
\]
\begin{enumerate}
\item Montrer que $\vect{EF}=\vect{FG}$. Que peut-on en déduire ?
\item Le point $G$ est-il le milieu du segment $[FH]$ ? Justifier la réponse.
\end{enumerate}

\exo

Dans un repère, on considère le point $A\couple{-1}{3}$, et le vecteur $\vecu\couple{2}{-5}$.

Le point $E\couple{x_E}{y_E}$, est tel que :
\[\vect{AE}=\vecu.\]
\begin{enumerate}
\item Justifier que l'on a : $\left\{\setlength\arraycolsep{2pt}\begin{array}{rcl}
x_E+1 & = & 2 \\
y_E-3 & = & -5
\end{array}\right.$
\item En déduire les coordonnées de $E$.
\end{enumerate}

\exo

Dans un repère, on considère les points :
\[
A\couple{2}{-5}
\qpvq
B\couple{2}{1}
\qpvq
C\couple{5}{-1}.
\]
Calculer les coordonnées du vecteur $\vect{BC}$, puis celles du point $M$ défini par $\vect{AM}=\vect{BC}$.

\exo

Dans un repère, on considère les points :
\[
A\couple{2}{-4}
\qpvq
B\couple{-2}{1}
\qpvq
C\couple{-5}{3}.
\]
Dans chaque cas, calculer les coordonnées du point $E$ :
\begin{enumerate}
\item $E$ est l'image de $A$ par la translation de vecteur $\vect{BC}$ ;
\item $E$ est le milieu de $[BC]$ ;
\item $E$ est le symétrique de $B$ par rapport à $A$.
\end{enumerate}

\exo

Dans un repère orthonormé, on considère les points :
\[
A\couple{5}{3}
\qpvq
B\couple{2}{-1}
\qpvq
C\couple{0}{3}.
\]
\begin{enumerate}
\item Calculer les distances $AB$, $BC$ et $AC$.
\item Quelle est la nature du triangle $ABC$ ?
\end{enumerate}

\exo

Dans un repère orthonormé, on considère les points :
\[
A\couple{1}{2}
\qpvq
B\couple{4}{-1}
\qpvq
C\couple{3}{1}.
\]
Montrer que le point $C$ appartient à la médiatrice du segment $[AB]$.

\exo

Dans un repère orthonormé, on considère les points :
\[
M\couple{-1}{-2}
\qpvq
N\couple{1}{4}
\qpvq
P\couple{4}{3}.
\]
Montrer que le triangle $MNP$ est rectangle en $N$.

\exo

Dans un repère orthonormé, on considère les points :
\[
B\couple{1}{-2}
\qpvq
C\couple{-3}{4}
\qpvq
E\couple{0}{6}.
\]
Montrer que les droites $(BC)$ et $(EC)$ sont perpendiculaires.

\exo

Dans un repère, on considère les points :
\[
A\couple{3}{-2}
\qpvq
B\couple{5}{4}
\qpvq
C\couple{1}{8}.
\]
On note $M$, $N$ et $P$ les milieux respectifs des segments $[AB]$, $[BC]$ et $[CA]$.
\begin{enumerate}
\item Calculer les coordonnées des points $M$, $N$ et $P$.
\item Montrer que le quadrilatère $MNPA$ est un parallélogramme.
\end{enumerate}

\exo

Dans un repère, on considère les points :
\[
A\couple{2}{-2}
\epve
B\couple{10}{2}
\epve
M\couple{8}{6}
\epve
N\couple{0}{2}.
\]
\begin{enumerate}
\item Montrer que le quadrilatère $ABMN$ est un parallélogramme.
\item Calculer $AM$ et $BN$. Que peut-on en déduire ?
\item Calculer l'aire de $ABMN$.
\end{enumerate}

\exo

Dans un repère, on considère les points :
\[
E\couple{1}{-3}
\epve
F\couple{9}{1}
\epve
G\couple{13}{-7}
\epve
H\couple{5}{-11}.
\]
Calculer l'aire du quadrilatère $EFGH$.

\exo

Soit $ABCD$ un quadrilatère. On note $I$, $J$, $K$ et $L$ les milieux des segments $[AB]$, $[BC]$, $[CD]$ et $[DA]$ :
\begin{center}
\def\xmin{0}
\def\xmax{9}
\def\ymin{-0.5}
\def\ymax{5.5}
\def\xunite{0.8cm}
\def\yunite{0.8cm}
\psset{unit=\xunite, xunit=\xunite, yunit=\yunite, algebraic=true, dimen=middle, dotstyle=*, dotsize=3pt 0, linewidth=0.8pt, arrowsize=3pt 2, arrowinset=0.25, comma=true}
\begin{pspicture*}(\xmin,\ymin)(\xmax,\ymax)
%\psgrid[xunit=\xunite, yunit=\yunite, subgriddiv=0, gridlabels=0, gridcolor=lightgray](0,0)(\xmin,\ymin)(\xmax,\ymax)
%\psaxes[labelFontSize=\scriptstyle, xAxis=true, yAxis=true, Dx=1, Dy=1, ticksize=-2pt 0, subticks=1]{->}(0,0)(\xmin,\ymin)(\xmax,\ymax)[{\footnotesize $x$},135][{\footnotesize $y$},-45]
\pspolygon(1,1)(7,0)(8,2)(4,5)
\pspolygon[fillcolor=black,fillstyle=solid,opacity=0.1](4,0.5)(7.5,1)(6,3.5)(2.5,3)
\begin{footnotesize}
\uput{5pt}[180](1,1){$A$}
\uput{5pt}[270](7,0){$B$}
\uput{5pt}[000](8,2){$C$}
\uput{5pt}[090](4,5){$D$}
\uput{5pt}[250](4,0.5){$I$}
\uput{5pt}[340](7.5,1){$J$}
\uput{5pt}[070](6,3.5){$K$}
\uput{5pt}[160](2.5,3){$L$}
\end{footnotesize}
\end{pspicture*}
\end{center}
\begin{enumerate}
\item 
\begin{enumerate}
\item Déterminer le réel $k$ tel que $\vect{AB}=k\vect{IB}$.
\item Exprimer de même $\vect{BC}$ en fonction de $\vect{BJ}$.
\item En déduire que $\vect{AC}=2\vect{IJ}$.
\end{enumerate}
\item Démontrer que le quadrilatère $IJKL$ est un parallélogramme.
\end{enumerate}

